% SOURCE_AUTHORITY: sga5_fr_workpass.tex, lines 5610--5629 (LF-normalized unit).
% SOURCE_SHA256: 791F4EFFC5E02832D5D77ED03518C8156D6F07E4C8238B03545DB93D883FBB28
\subsection*{6.3}
Se $t:T\to S$ é um $S$-esquema e $A$ uma $\Lambda$-álgebra, poremos
\begin{equation}
t^!A = KA_T \qquad(\in\Ob D^+({}_AT_A)).
\tag{6.3.1}
\end{equation}
Para $A=\Lambda$, escreveremos $K_T$ em vez de $K\Lambda_T$. De 6.2.2 resulta que se tem um isomorfismo canônico em $D^+({}_AT_A)$
\begin{equation}
K_T\lten A\xrightarrow{\sim}KA_T.
\tag{6.3.2}
\end{equation}
Denotaremos por $D_A$ o funtor
\[
\uRHom_A(-,KA_T).
\]
Este envia $D_{\mathrm{ctf}}({}_AT)$ a $D_{\mathrm{ctf}}(T_A)$ e vice-versa, e a demonstração de (SGA $4^{1/2}$ Th. Finitude 4) mostra que, para $L\in\Ob D_{\mathrm{ctf}}({}_AT)$, a flecha canônica
\[
L\longrightarrow D_AD_AL
\]
é um isomorfismo. Para $A=\Lambda$, escreveremos $D$ em vez de $D_A$.
